% Generated draft: requires mathematical and source review.
\documentclass[11pt]{article}
\usepackage[margin=25mm]{geometry}
\usepackage{fontspec}
\setmainfont{Latin Modern Roman}
\usepackage{amsmath,amssymb,mathtools,unicode-math}
\setmathfont{Latin Modern Math}
\usepackage{enumitem,needspace,xurl,hyperref}
\setlength{\parindent}{0pt}
\setlength{\parskip}{5pt}
\setlength{\emergencystretch}{3em}
\newcommand{\N}{\mathbb{N}}
\newcommand{\Z}{\mathbb{Z}}
\newcommand{\Q}{\mathbb{Q}}
\newcommand{\R}{\mathbb{R}}
\newcommand{\C}{\mathbb{C}}
\newcommand{\F}{\mathbb{F}}
\newcommand{\A}{\mathbb{A}}
\newcommand{\PP}{\mathbb{P}}
\newcommand{\E}{\mathbb{E}}
\newcommand{\eps}{\varepsilon}
\newcommand{\dd}{\,\mathrm{d}}
\begin{document}
\section*{UnsolvedMath ID: 4200008}
\textbf{Source problem code:} catalogue entry
\subsection{Definitions and mathematical statement}
Let \(\Omega\subseteq\R^d\) be the fixed convex shape of a rigid container. Translation-only motion means that its occupied region is \(q(t)+\Omega\), with velocity \(V(t)=\dot q(t)\), and no time-dependent rotation. A free gas has ballistic motion between collisions, with no interparticle forces; its exchanges of momentum with the wall couple its evolution to \(q\). For a continuum gas of phase-space density \(f_t(x,v)\), the formal total momentum is
\[
P=M V(t)+\int_{q(t)+\Omega}\int_{\R^d}v f_t(x,v)\,\mathrm dv\,\mathrm dx,
\]
where \(M>0\) is the container mass and \(f_t\) is normalized as mass density. Conservation means \(P\) is independent of time. This describes a possible continuum framework, not collision laws or existence assumptions supplied by the extracted record.

To compare densities on a fixed spatial domain, let \(\mu_t\) be the gas phase-space measure expressed in translated coordinates \((y,v)=(x-q(t),v)\in\Omega\times\R^d\). Weak convergence to a limiting measure \(\mu_\infty\) would mean
\[
\forall\varphi\in C_b(\Omega\times\R^d),\qquad
\int\varphi\,\mathrm d\mu_t\longrightarrow\int\varphi\,\mathrm d\mu_\infty
\quad(t\to\infty).
\]
The question asks whether this occurs with an equilibrium limit and with \(V(t)\to0\). However, the class of initial data, the equilibrium condition, and the reference frame are not specified. In particular, coming to rest is frame-dependent; a zero-total-momentum or other chosen frame cannot be silently assumed. Neither a Maxwellian limit nor convergence of the container position is stated by the source.

The source also does not distinguish finitely many particles from a continuum density model, or give the collision rule, dimension, boundary regularity, or masses needed to define the coupled dynamics. The supplied research summary reports that existence of the dynamics was solved by Knill, while weak convergence to equilibrium and rest of the container remain conjectural with partial or related results. That existence assessment does not settle the asymptotic question.
\subsection{Short English statement}
Does a collisionless gas exchanging momentum with a convex translating container converge weakly to equilibrium while the container velocity tends to zero? The dynamical model, admissible data, equilibrium, and reference frame must be specified.
\subsection{Sources}
\begin{itemize}
\item Local UnsolvedMath registry, permanent ID \(4200008\), exact source code \texttt{catalogue entry}, ``Free gas in a moving container.''
\item Oliver Knill, \emph{Open problems in Hamiltonian dynamics} (2000), item 8, ``FREE GAS IN MOVING CONTAINER''; \url{https://people.math.harvard.edu/~knill/seminars/intr/}. HTML extraction; recorded access 2026-07-29.
\item Supplied research summary attributes existence of the dynamics to Knill and describes the convergence question as open, with partial/related results. It supplies no complete publication reference or precise existence theorem.
\end{itemize}
\end{document}
